% !TEX root = ../discretemath.tex

\section{Матричные игры}

\subsection{Постановка задачи}

Рассматривается игра двух игроков с нулевой суммой. Задаётся \emph{матрица выигрышей}
\[
A = \begin{pmatrix} a_{11} & \cdots & a_{1n} \\ \vdots & & \vdots \\ a_{m1} & \cdots & a_{mn} \end{pmatrix}
\in \mathbb{R}^{m \times n}.
\]
Правила игры:
\begin{enumerate}
	\item 1-й игрок \emph{независимо} выбирает строку $i \in \{1, \ldots, m\}$.
	\item 2-й игрок \emph{независимо} выбирает столбец $j \in \{1, \ldots, n\}$.
	\item 2-й игрок выплачивает 1-му сумму $a_{ij}$.
\end{enumerate}
Таким образом, 1-й игрок стремится \emph{максимизировать} выигрыш, а 2-й --- \emph{минимизировать} его.

\subsection{Чистые стратегии. Цена игры и седловая точка}

\begin{Definition}{Нижняя и верхняя цены игры}{}
	\emph{Нижняя цена игры} --- наилучшее, что может гарантировать 1-й игрок:
	\[
	\alpha = \max_{i}\,\min_{j}\, a_{ij}.
	\]
	\emph{Верхняя цена игры} --- наилучшее, что может гарантировать 2-й игрок:
	\[
	\beta = \min_{j}\,\max_{i}\, a_{ij}.
	\]
\end{Definition}

Иными словами, выбирая строку $i^* = \arg\max_i \min_j a_{ij}$, 1-й игрок \emph{обеспечивает} себе выигрыш не менее~$\alpha$ при любом ходе соперника. Аналогично, 2-й игрок, выбирая $j^* = \arg\min_j \max_i a_{ij}$, \emph{ограничивает} свои потери сверху величиной~$\beta$.

\begin{Theorem}{Для любой матрицы $\alpha \leq \beta$}{}
	Для любой матрицы $A \in \mathbb{R}^{m \times n}$ выполнено $\alpha \leq \beta$.
\end{Theorem}

\begin{proof}
Для любых фиксированных $i, j$ имеем
$\min_{j'} a_{ij'} \leq a_{ij} \leq \max_{i'} a_{i'j}$.
Беря максимум по $i$ в левой части и минимум по $j$ в правой:
$\alpha = \max_i \min_{j'} a_{ij'} \leq \min_j \max_{i'} a_{i'j} = \beta$. 
\end{proof}

\begin{Definition}{Седловая точка}{}
	Если $\alpha = \beta$, то говорят, что у игры есть \emph{цена} $\upsilon = \alpha = \beta$ и \emph{седловая точка}~$(i^*, j^*)$, удовлетворяющая условию
	\[
	a_{ij^*} \leq a_{i^*j^*} \leq a_{i^*j} \quad \text{для всех } i \in \{1,\ldots,m\},\; j \in \{1,\ldots,n\}.
	\]
	Это означает, что $a_{i^*j^*}$ одновременно является минимумом в строке $i^*$ и максимумом в столбце $j^*$.
\end{Definition}

\begin{Remark}{}{}
	Если заменить матрицу $A$ на $-A^{\top}$, роли игроков меняются местами.
\end{Remark}

\begin{Statement}{О единственности значения в седловых точках}{}
	Если у матрицы $A$ есть два седла с индексами $(i_1, j_1)$ и $(i_2, j_2)$, то все четыре элемента
	$a_{i_1 j_1},\, a_{i_1 j_2},\, a_{i_2 j_1},\, a_{i_2 j_2}$ равны $\upsilon$.
\end{Statement}

\begin{proof}
 Из условия седла в $(i_1, j_1)$: $a_{i_2 j_1} \leq a_{i_1 j_1}$ и $a_{i_1 j_1} \leq a_{i_1 j_2}$.
Из условия седла в $(i_2, j_2)$: $a_{i_1 j_2} \leq a_{i_2 j_2}$ и $a_{i_2 j_2} \leq a_{i_2 j_1}$.
Тогда
\[
a_{i_1 j_1} \leq a_{i_1 j_2} \leq a_{i_2 j_2} \leq a_{i_2 j_1} \leq a_{i_1 j_1},
\]
откуда все четыре значения совпадают.
\end{proof}

\begin{Example}{Матрица $3\times 3$ с седловой точкой}{}
	Рассмотрим матрицу
	\[
	A = \begin{pmatrix} 2 & 4 & 7 \\ 6 & \mathbf{5} & 6 \\ 3 & 4 & 1 \end{pmatrix}.
	\]
	Минимумы по строкам: $2,\; 5,\; 1$, поэтому $\alpha = \max(2, 5, 1) = 5$.
	Максимумы по столбцам: $6,\; 5,\; 7$, поэтому $\beta = \min(6, 5, 7) = 5$.
	Так как $\alpha = \beta = 5$, цена игры равна $5$. Седловая точка --- элемент $a_{22} = 5$.
\end{Example}

\subsection{Смешанные стратегии}

\subsubsection*{Определения}

При чистых стратегиях каждый игрок может гарантировать лишь $\alpha \leq \upsilon \leq \beta$. Чтобы всегда достигать точной цены $\upsilon$, вводят \emph{смешанные} (рандомизированные) стратегии.

\begin{Definition}{Смешанные стратегии и ожидаемый выигрыш}{}
	\emph{Смешанная стратегия 1-го игрока} --- вектор вероятностей
	\[
	\vec{p} = (p_1, \ldots, p_m), \quad p_i \geq 0, \quad \sum_{i=1}^m p_i = 1.
	\]
	\emph{Смешанная стратегия 2-го игрока} --- вектор вероятностей
	\[
	\vec{q} = (q_1, \ldots, q_n), \quad q_j \geq 0, \quad \sum_{j=1}^n q_j = 1.
	\]
	\emph{Ожидаемый выигрыш 1-го игрока} при смешанных стратегиях:
	\[
	E(\vec{p},\,\vec{q}) = \sum_{i=1}^{m}\sum_{j=1}^{n} a_{ij}\, p_i\, q_j.
	\]
\end{Definition}

\begin{Definition}{Оптимальная пара стратегий}{}
	Пара $(\vec{p}^{\,*},\, \vec{q}^{\,*})$ называется \emph{оптимальной} (или \emph{равновесием в смешанных стратегиях}), если для всех смешанных стратегий $\vec{p}$ и $\vec{q}$:
	\[
	E(\vec{p},\, \vec{q}^{\,*}) \leq E(\vec{p}^{\,*},\, \vec{q}^{\,*}) \leq E(\vec{p}^{\,*},\, \vec{q}).
	\]
	Число $\upsilon = E(\vec{p}^{\,*},\, \vec{q}^{\,*})$ называется \emph{ценой игры}.
\end{Definition}

\subsubsection*{Свойства множества оптимальных пар}

\begin{Statement}{О выпуклости множества оптимумов}{}
	Пусть $(\vec{p}^{\,*},\, \vec{q}^{\,*})$ и $(\vec{p}^{\,\prime},\, \vec{q}^{\,\prime})$ --- две оптимальные пары. Тогда:
	\begin{enumerate}
		\item $E(\vec{p}^{\,*},\, \vec{q}^{\,*}) = E(\vec{p}^{\,\prime},\, \vec{q}^{\,\prime}) = \upsilon$ (все оптимальные пары дают одинаковую цену).
		\item Для любых $\alpha,\, \beta \in [0,1]$ пара
		\[
		\bigl(\alpha\vec{p}^{\,*} + (1-\alpha)\vec{p}^{\,\prime},\;\; \beta\vec{q}^{\,*} + (1-\beta)\vec{q}^{\,\prime}\bigr)
		\]
		также является оптимальной.
	\end{enumerate}
\end{Statement}

\begin{proof}
Поскольку $(\vec{p}^{\,*},\, \vec{q}^{\,*})$ оптимальна:
\[
E(\vec{p}^{\,*},\, \vec{q}^{\,*}) \leq E(\vec{p}^{\,*},\, \vec{q}^{\,\prime}) \leq E(\vec{p}^{\,\prime},\, \vec{q}^{\,\prime}).
\]
Аналогично в обратную сторону, откуда $E(\vec{p}^{\,*},\, \vec{q}^{\,*}) = E(\vec{p}^{\,\prime},\, \vec{q}^{\,\prime}) =: \upsilon$.

Проверим пункт 2. Для произвольного $\vec{q}$ и $\alpha\vec{p}^{\,*} + (1-\alpha)\vec{p}^{\,\prime}$ в силу линейности:
\begin{align*}
	E\!\left(\vec{p},\; \beta\vec{q}^{\,*} + (1-\beta)\vec{q}^{\,\prime}\right)
	&= \beta\,E(\vec{p},\,\vec{q}^{\,*}) + (1-\beta)\,E(\vec{p},\,\vec{q}^{\,\prime})
	\leq \beta\,\upsilon + (1-\beta)\,\upsilon = \upsilon, \\[4pt]
	E\!\left(\alpha\vec{p}^{\,*} + (1-\alpha)\vec{p}^{\,\prime},\; \vec{q}\right)
	&= \alpha\,E(\vec{p}^{\,*},\,\vec{q}) + (1-\alpha)\,E(\vec{p}^{\,\prime},\,\vec{q})
	\geq \alpha\,\upsilon + (1-\alpha)\,\upsilon = \upsilon.
\end{align*}
Следовательно, выпуклая комбинация удовлетворяет определению оптимальной пары.
\end{proof}

\subsubsection*{Пример нахождения оптимальных стратегий}

\begin{Example}{Матрица $2\times 2$}{}
	Пусть
	\[
	A = \begin{pmatrix} -1 & 3 \\ 3 & -5 \end{pmatrix}.
	\]
	Смешанные стратегии: 1-й игрок --- $(p,\,1-p)$, 2-й --- $(q,\,1-q)$. Ожидаемый выигрыш:
	\begin{align*}
		E(p,q) &= (-1)\cdot pq + 3\cdot p(1-q) + 3\cdot(1-p)q + (-5)\cdot(1-p)(1-q) \\
		&= -5 + 8p + 8q - 12pq.
	\end{align*}
	
	Найдём оптимальную пару. При фиксированном $p$
	\[
	E(p,q)=(-5+8p)+(8-12p)q.
	\]
	Если $p^* \neq \tfrac23$, то коэффициент при $q$ ненулевой, и минимум по $q$ достигается в одной из чистых стратегий: $q^*=0$ или $q^*=1$.
	Но тогда оптимальный ответ 1-го игрока тоже будет чистым:
	\[
	E(p,0)=-5+8p \quad \Rightarrow \quad p^*=1,
	\qquad
	E(p,1)=3-4p \quad \Rightarrow \quad p^*=0,
	\]
	что невозможно, поскольку при $p^*=1$ минимум достигается при $q^*=1$, а при $p^*=0$ --- при $q^*=0$.
	Следовательно,
	\[
	p^*=\tfrac23.
	\]
	Аналогично,
	\[
	q^*=\tfrac23.
	\]
	
	Проверка:
	\[
	E\!\left(p,\tfrac23\right)=\tfrac13 \quad \text{для всех } p,
	\qquad
	E\!\left(\tfrac23,q\right)=\tfrac13 \quad \text{для всех } q.
	\]
	Значит,
	\[
	\vec p^{\,*}=\left(\tfrac23,\tfrac13\right), \qquad
	\vec q^{\,*}=\left(\tfrac23,\tfrac13\right), \qquad
	\upsilon=\tfrac13.
	\]
\end{Example}

\subsubsection*{Теорема существования и связь с линейным программированием}

\begin{Theorem}{О существовании оптимальной смешанной стратегии}{}
	В любой матричной игре $A \in \mathbb{R}^{m \times n}$ существует оптимальная пара смешанных стратегий $(\vec{p}^{\,*},\, \vec{q}^{\,*})$.
\end{Theorem}

Ниже изложена идея доказательства через сведение к линейному программированию. Без ограничения общности считаем $a_{ij} > 0$ (к любой матрице можно прибавить константу, что не меняет оптимальных стратегий). Тогда цена игры $\upsilon > 0$.

Если 2-й игрок «раскрывает» чистую стратегию $j$ (то есть играет $\vec{q} = e_j$), условие оптимальности для 1-го игрока принимает вид:
\[
\sum_{i=1}^m p_i^*\, a_{ij} \geq \upsilon \quad \forall\, j = 1,\ldots,n.
\]
Кроме того, $\upsilon = \max_{\vec{p}}\, \min_{j}\, \sum_i p_i\, a_{ij}$ --- это наибольшее число, которое 1-й игрок может гарантировать. После замены $x_i = p_i^*/\upsilon > 0$ задача нахождения оптимальной стратегии 1-го игрока сводится к следующей задаче линейного программирования (\textbf{прямая задача}):
\[
\begin{cases}
	\displaystyle\sum_{i=1}^m x_i\, a_{ij} \geq 1, \quad j = 1,\ldots, n, \\[4pt]
	x_1, \ldots, x_m \geq 0, \\[2pt]
	\displaystyle\sum_{i=1}^m x_i \;\longrightarrow\; \min,
\end{cases}
\]
при этом $p_i^* = x_i / \sum_k x_k$ и $\sum_i x_i = 1/\upsilon$ (минимизация $\sum x_i$ равносильна максимизации $\upsilon$).

\textbf{Двойственная задача} (стратегия 2-го игрока, замена $y_j = q_j^*/\upsilon$):
\[
\begin{cases}
	\displaystyle\sum_{j=1}^n a_{ij}\, y_j \leq 1, \quad i = 1,\ldots, m, \\[4pt]
	y_1, \ldots, y_n \geq 0, \\[2pt]
	\displaystyle\sum_{j=1}^n y_j \;\longrightarrow\; \max = \tfrac{1}{\upsilon}.
\end{cases}
\]
По теореме двойственности в ЛП обе задачи имеют решения, что и гарантирует существование оптимальной пары.

\subsection{Биматричные игры}

\subsubsection*{Постановка задачи}

\begin{Definition}{Биматричная игра}{}
	\emph{Биматричная игра} задаётся парой матриц $M_1,\, M_2 \in \mathbb{R}^{m \times n}$:
	\begin{itemize}
		\item 1-й игрок выбирает строку $i \in \{1,\ldots, m\}$,
		\item 2-й игрок выбирает столбец $j \in \{1,\ldots, n\}$,
		\item выигрыш 1-го равен $M_1(i, j)$, выигрыш 2-го равен $M_2(i, j)$.
	\end{itemize}
	При смешанных стратегиях $\vec{p}$, $\vec{q}$ ожидаемые выигрыши:
	\[
	M_1(\vec{p},\,\vec{q}) = \sum_{i,j} M_1(i,j)\, p_i\, q_j, \qquad
	M_2(\vec{p},\,\vec{q}) = \sum_{i,j} M_2(i,j)\, p_i\, q_j.
	\]
\end{Definition}

\begin{Remark}{}{}
	Матричная игра $A$ является частным случаем биматричной: $(M_1, M_2) = (A,\,-A)$.
\end{Remark}

\subsubsection*{Равновесие Нэша}

\begin{Definition}{Равновесие Нэша}{}
	Пара смешанных стратегий $(\vec{p}^{\,*},\, \vec{q}^{\,*})$ называется \emph{равновесием Нэша}, если ни один из игроков не может увеличить свой выигрыш, отклонившись в одностороннем порядке:
	\[
	M_1(\vec{p},\, \vec{q}^{\,*}) \leq M_1(\vec{p}^{\,*},\, \vec{q}^{\,*}) \quad \forall\, \vec{p},
	\]
	\[
	M_2(\vec{p}^{\,*},\, \vec{q}) \leq M_2(\vec{p}^{\,*},\, \vec{q}^{\,*}) \quad \forall\, \vec{q}.
	\]
\end{Definition}

\begin{Example}{Дилемма заключённого}{}
	Пусть $M_1, M_2$ — матрицы выигрышей:
	\[
	M_1 = \begin{pmatrix} -8 & 0 \\ -10 & -1 \end{pmatrix}, \qquad
	M_2 = \begin{pmatrix} -8 & -10 \\ 0 & -1 \end{pmatrix}.
	\]
	Стратегия $i=1$ — «признаться», $i=2$ — «молчать».
	
	Для каждого игрока «признаться» является \textbf{доминирующей стратегией}, так как $-8 > -10$ (если другой признался) и $0 > -1$ (если другой молчит). Таким образом, единственное равновесие Нэша — стратегия $(1, 1)$, где оба получают по 8 лет.
	
	\textbf{Суть дилеммы:} Исход $(-1, -1)$ при совместном молчании был бы коллективно выгоднее, однако он неустойчив. Рациональный эгоизм заставляет каждого игрока выбрать предательство, что приводит к худшему общему результату.
\end{Example}

\begin{Example}{Семейный спор}{}
	Пусть
	\[
	M_1 = \begin{pmatrix} 1 & 0 \\ 0 & 4 \end{pmatrix}, \qquad
	M_2 = \begin{pmatrix} 4 & 0 \\ 0 & 1 \end{pmatrix}.
	\]
	Смешанные стратегии: $\vec{p} = (p,\, 1-p)$, $\vec{q} = (q,\, 1-q)$.
	
	\begin{enumerate}
		\item $\vec{p}^{\,*} = (1, 0),\; \vec{q}^{\,*} = (1, 0)$ --- равновесие Нэша с выигрышами $(1, 4)$.
		\item $\vec{p}^{\,*} = (0, 1),\; \vec{q}^{\,*} = (0, 1)$ --- равновесие Нэша с выигрышами $(4, 1)$.
		\item Поиск внутреннего равновесия ($0 < p^*, q^* < 1$):
		\[
		M_1(p, q^*) = p(5q^*-4) + 4 - 4q^* \;\Longrightarrow\; q^* = 4/5.
		\]
		Аналогично для $M_2$: $p^* = 1/5$.
		
		Выигрыши в смешанном равновесии $\vec{p}^{\,*} = (1/5, 4/5), \vec{q}^{\,*} = (4/5, 1/5)$:
		\[
		v_1 = v_2 = \tfrac{1}{5}\cdot\tfrac{4}{5}\cdot 1 + \tfrac{4}{5}\cdot\tfrac{1}{5}\cdot 4 = \tfrac{4}{5}.
		\]
	\end{enumerate}
	\textbf{Вывод:} Смешанное равновесие неэффективно, так как выигрыш в нем ($0.8$) меньше минимально возможного выигрыша в любом из чистых равновесий ($1$). Отсутствие координации ведет к потерям для обоих игроков.
\end{Example}

\subsubsection*{Теорема о существовании равновесия Нэша}

\begin{Theorem}{О существовании равновесия Нэша}{}
	В любой биматричной игре существует хотя бы одно равновесие Нэша.
\end{Theorem}

Доказательство опирается на теорему Какутани о неподвижной точке, которую сформулируем ниже.

\begin{Theorem}{Какутани о неподвижной точке}{}
	Пусть $S \subset \mathbb{R}^n$ --- непустое компактное выпуклое множество, и $F \colon S \to 2^S$ --- многозначное отображение такое, что:
	\begin{enumerate}
		\item $F(x) \neq \emptyset$ для всех $x \in S$,
		\item $F(x)$ выпукло для всех $x \in S$,
		\item $F$ \emph{замкнуто}: если $x_k \to x$ и $y_k \to y$ с $y_k \in F(x_k)$, то $y \in F(x)$.
	\end{enumerate}
	Тогда существует неподвижная точка: $\exists\, x^* \in S$ такая, что $x^* \in F(x^*)$.
\end{Theorem}

\begin{Remark}{}{}
	Теорема Какутани обобщает классическую теорему Брауэра о неподвижной точке (непрерывное отображение замкнутого шара в себя имеет неподвижную точку) на случай многозначных замкнутых отображений.
\end{Remark}

\textbf{Доказательство существования равновесия Нэша.}
Рассмотрим компактное выпуклое множество всех пар смешанных стратегий:
\[
S = \left\{(\vec{p},\,\vec{q}) :\ p_i \geq 0,\; \sum_{i=1}^m p_i = 1,\; q_j \geq 0,\; \sum_{j=1}^n q_j = 1\right\} \subset \mathbb{R}^{m+n}.
\]
Определим многозначное отображение $F \colon S \to 2^S$:
\[
F(\vec{p},\,\vec{q}) = \left\{(\vec{p}{\,}^{\,\prime},\,\vec{q}{\,}^{\,\prime}) \in S :
\begin{array}{l}
	M_1(\vec{p}{\,}^{\,\prime},\,\vec{q}) = \max_{\vec{x}}\, M_1(\vec{x},\,\vec{q}), \\[2pt]
	M_2(\vec{p},\,\vec{q}{\,}^{\,\prime}) = \max_{\vec{y}}\, M_2(\vec{p},\,\vec{y})
\end{array}
\right\}.
\]
Проверим условия теоремы Какутани:
\begin{enumerate}
	\item \textbf{Непустота.} Максимум непрерывной функции на компакте $S$ всегда достигается, поэтому $F(\vec{p},\,\vec{q}) \neq \emptyset$.
	\item \textbf{Выпуклость.} Функции $M_1(\,\cdot\,,\vec{q})$ и $M_2(\vec{p},\,\cdot\,)$ линейны по первому и второму аргументу соответственно. Поэтому множество максимизаторов выпукло, и любая линейная комбинация элементов $F(\vec{p},\,\vec{q})$ снова лежит в $F(\vec{p},\,\vec{q})$.
	\item \textbf{Замкнутость.} Если $(\vec{p}{\,}^{(k)},\, \vec{q}{\,}^{(k)}) \to (\vec{p},\, \vec{q})$ и $(\vec{p}{\,}'^{(k)},\, \vec{q}{\,}'^{(k)}) \to (\vec{p}{\,}',\, \vec{q}{\,}')$, причём $(\vec{p}{\,}'^{(k)},\, \vec{q}{\,}'^{(k)}) \in F(\vec{p}{\,}^{(k)},\, \vec{q}{\,}^{(k)})$, то по непрерывности $M_1$ и $M_2$ получаем $(\vec{p}{\,}',\, \vec{q}{\,}') \in F(\vec{p},\, \vec{q})$.
\end{enumerate}
По теореме Какутани существует $(\vec{p}^{\,*},\,\vec{q}^{\,*}) \in F(\vec{p}^{\,*},\,\vec{q}^{\,*})$, что в точности означает условие равновесия Нэша:
\[
M_1(\vec{p}^{\,*},\,\vec{q}^{\,*}) = \max_{\vec{x}}\, M_1(\vec{x},\,\vec{q}^{\,*}), \qquad
M_2(\vec{p}^{\,*},\,\vec{q}^{\,*}) = \max_{\vec{y}}\, M_2(\vec{p}^{\,*},\,\vec{y}). \quad \square
\]